% GENERATED FILE. Edit canonical lessons, then run npm run generate:books.
% canonical-source-hash: fnv1a64:0f9658a0088c0b7d
% canonical-lessons: TA-C30-maalai-vanakkam

\chapter{Good Evening}
\label{ch:ta-good-evening}

\begin{chapteropening}
\textbf{By the end of this chapter:} \emph{I can greet someone with \ta{மாலை} \ta{வணக்கம்} from late afternoon through nightfall, built on the same two-piece pattern as the morning greeting.}

Greet with mālai vaṇakkam, assembling mālai plus vaṇakkam, and note the pattern was re-checked rather than assumed to carry over.
\end{chapteropening}

\section[\ta{மாலை} \ta{வணக்கம்}]{\ta{மாலை} \ta{வணக்கம்} — \textquotedblleft{}evening greetings\textquotedblright{}}

\label{lesson:TA-C30-maalai-vanakkam}



\begin{quote}
Last chapter's morning greeting followed its own pattern. This lesson checks whether the evening greeting follows the same one — rather than assuming it automatically does.
\end{quote}

\subsection*{You'll want to know: \ta{மாலை} \ta{வணக்கம்} — the same compositional pattern, independently confirmed}
\textbf{\ta{மாலை} \ta{வணக்கம்}} (\textbf{mālai vaṇakkam}) — \textquotedblleft{}\textbf{good evening}\textquotedblright{} — is confirmed, by two independently-fetched sources, as a real, standard Tamil evening greeting, used roughly from \textbf{afternoon through nightfall} (the two sources give slightly different exact windows — \textquotedblleft{}after noon through evening\textquotedblright{} versus \textquotedblleft{}late afternoon until nightfall\textquotedblright{} — a small, unsurprising looseness, matching the noon/afternoon boundary-blur already seen in TA-C28). It's built exactly the same way as TA-C29's morning greeting: \textbf{\ta{மாலை}} (already met, Chapter 27, \textquotedblleft{}\textbf{evening}\textquotedblright{}) plus \textbf{\ta{வணக்கம்}} (already met, Chapter 1, Tamil's general greeting). Literally: \textquotedblleft{}\textbf{evening salutations}.\textquotedblright{}

\begin{culture}
Here's the honest methodological point: just because TA-C29's morning greeting followed the \textquotedblleft{}time-word + vaṇakkam\textquotedblright{} pattern doesn't mean this evening greeting was assumed to follow it automatically. This arc has repeatedly found that greeting formation doesn't reliably generalize, even within the very same language — Malayalam's own morning greeting broke from the śubha+time-word pattern its afternoon and evening greetings both share, so \textquotedblleft{}the pattern probably repeats\textquotedblright{} was never a safe assumption on its own. So this greeting was independently verified via direct fetch before confirming the pattern really does repeat here — and it does, cleanly: no śubha-equivalent word appears in Tamil's evening greeting either, same as the morning one.
\end{culture}

\subsection*{Your turn}
Say these aloud:

\begin{itemize}
\raggedright
  \item \textquotedblleft{}mālai vaṇakkam\textquotedblright{} — \textquotedblleft{}good evening,\textquotedblright{} literally \textquotedblleft{}evening greetings\textquotedblright{}
  \item the two pieces you already know — mālai (Chapter 27) + vaṇakkam (Chapter 1)
  \item the honest methodological point — this pattern was re-verified, not assumed to carry over automatically
\end{itemize}

\begin{tcolorbox}[breakable,colback=teal!4,colframe=teal!35!black,title={Before you move on}]
What two already-known words does \ta{மாலை} \ta{வணக்கம்} build from? (\ta{மாலை}, \textquotedblleft{}evening\textquotedblright{} + \ta{வணக்கம்}, \textquotedblleft{}greetings.\textquotedblright{}) Was it assumed that the evening greeting would follow the same pattern as the morning one, or was it independently checked? (\textbf{Independently checked} — via direct fetch, not assumed, since this arc has learned that greeting patterns don't reliably generalize.) Does \ta{மாலை} \ta{வணக்கம்} use a śubha-equivalent word, like Hindi/Kannada/Telugu/Malayalam's evening greetings do? (\textbf{No} — same as Tamil's morning greeting, no śubha-equivalent word appears here either.)
\end{tcolorbox}
